\section{14.1 MuJoCo와 정책 학습 실습}
\begin{frame}{장난감 환경에서 실체 물리엔진으로}
  \begin{itemize}
    \item Week 13의 Gymnasium 기본 환경: 물리를 \textbf{간단한 근사식}으로
      계산 --- 빠르지만 접촉·마찰은 거칠게 흉내만 냄
    \item 실전 로봇 연구: 접촉(contact), 마찰 등 중요한 물리 현상은
      \textbf{훨씬 정교한 물리엔진}이 필요
  \end{itemize}
  \vskip0.6em
  \small
  \begin{tabular}{@{}ll@{}}
    \toprule
    도구 & 특징 \\
    \midrule
    \kb{MuJoCo} (2012, Todorov) & CPU에서 도는 \textbf{정밀} 물리엔진 ---
      접촉·마찰을 LCP로 매 스텝 해결 \\
    \kb{Isaac Sim} (NVIDIA) & GPU로 \textbf{수천 개} 시뮬레이션을
      동시에 --- 렌더링 포함 \\
    \bottomrule
  \end{tabular}
  \vskip0.7em
  \begin{block}{이번 장이 줄곧 묻는 질문}
    \kq{같은 강화학습 알고리즘 위에 시뮬레이터를 교체했을 때,
    무엇이 바뀌고 무엇이 그대로 남는가?}
  \end{block}
  정밀하고 빠른 물리는 어디까지 필요한가, 그 비용을 누가 지불하는가
  --- 실측 숫자까지 붙여 답한다.
\end{frame}
